\subsection{Splitting extensions}

DEFINITION 2. --- Let K be a field and \((f_{i})_{i\in I}\) a family of non-constant polynomials in K{[}X{]}. By a splitting extension of \((f_{i})_{i\in I}\) we understand any extension E of K with the following properties:

\begin{enumerate}
\def\labelenumi{\alph{enumi})}
\item
  For each \(i \in \mathbf{I}\) , the polynomial \(f_i\) splits in \(\mathbf{E}[\mathbf{X}]\) into a product of polynomials of degree 1.
\item
  For each \(i \in I\) let \(\mathbf{R}_i\) be the set of roots of \(f_i\) in \(E\) , then \(E = K\left(\bigcup_{i \in I} \mathbf{R}_i\right)\) .
\end{enumerate}

Sometimes the term « splitting field » is used instead of « splitting extension ».

Remarks. --- 1) For each \(i \in I\) let \(c_i\) be a non-zero element of K and let \(f_i' = c_i f_i\) . Then it is clear that every splitting extension for the family \((f_i)_{i \in I}\) is also a splitting extension for the family \((f_i')_{i \in I}\) and conversely. In particular, in studying splitting extensions we may limit ourselves to the case of monic polynomials.

\begin{enumerate}
\def\labelenumi{\arabic{enumi})}
\setcounter{enumi}{1}
\tightlist
\item
  Suppose that I is finite and put \(f = \prod_{i \in I} f_i\) . Using the uniqueness of the
\end{enumerate}

decomposition of a polynomial into irreducible factors in E{[}X{]} (IV, p.~13, Prop. 13), we can easily show that a splitting extension for the polynomial f is a splitting extension for the family \((f_{i})_{i \in I}\) , and conversely. In other words, the case of a finite family may be reduced to the case of a single polynomial.

\begin{enumerate}
\def\labelenumi{\arabic{enumi})}
\setcounter{enumi}{2}
\tightlist
\item
  Let \(f \in K[X]\) be a polynomial of degree \(\geqslant 1\) and let E be a splitting extension of f.~If \(x_{1}, \ldots, x_{n}\) are the roots of f in E, we thus have \(\mathrm{E} = \mathrm{K}(x_{1}, \ldots, x_{n})\) and \([E : K]\) is finite (V, p.~18, Th. 2); but it may happen that E is distinct from the subfields \(\mathrm{K}(x_{1}), \ldots, \mathrm{K}(x_{n})\) generated by a single root; this can happen even when f is irreducible \(^{1}\) . We note however that when f is irreducible, the fields \(\mathrm{K}(x_{i})\) all have the same degree n over K, and whenever E is equal to one of them, we have \([E : K] = n\) and hence \(\mathrm{E} = \mathrm{K}(x_{1}) = \cdots = \mathrm{K}(x_{n})\) .
\end{enumerate}

PROPOSITION 4. --- Let K be a field and \((f_{i})_{i\in I}\) a family of non-constant polynomials in K{[}X{]}; then there exists a splitting extension for the family \((f_{i})_{i\in I}\) .

We may take the polynomials \(f_{i}\) to be monic (Remark 1). Let \(i \in I\) and let the degree of \(f_{i}\) be \(d_{i}\) . By IV, p.~73, Prop. 5 there exists a commutative algebra \(A_{i}\) over \(K\) , not reduced to 0, and elements \(\xi_{i,1}, \ldots, \xi_{i,d_{i}}\) of \(A_{i}\) such that:

\begin{enumerate}
\def\labelenumi{\alph{enumi})}
\item
  the algebra \(\mathbf{A}_i\) is generated by \((\xi_{i,1},\dots,\xi_{i,d_i})\) ;
\item
  we have \(f_{i}(\mathbf{X}) = \prod_{k=1}^{d_{i}} (\mathbf{X} - \xi_{i,k})\) in \(\mathbf{A}_i[\mathbf{X}]\) .
\end{enumerate}

Let A be the tensor product of the family of algebras \((\mathbf{A}_{i})_{i\in I}\) and let \(\varphi_{i}\) be the canonical homomorphism of \(A_{i}\) into A (III, p.~470). The algebra A is then commutative and not reduced to 0; by Krull's theorem (I, p.~104), there exists thus a maximal ideal a in A and \(E = A/a\) is an extension of the field K.

Denote by \(\psi\) the canonical homomorphism of A into E and put \(x_{i,k} = \psi (\varphi_i(\xi_{i,k}))\) for \(i\in I\) and \(1\leqslant k\leqslant d_i\) . Since the algebra A is generated by \(\cup \varphi_i(A_i)\) , the extension E is generated by the family \((x_{i,k})\) . Further, we have \(i\in I\)

\(f_{i}(\mathbf{X}) = \prod_{k=1}^{d_{i}} (\mathbf{X} - x_{i,k})\) in E{[}X{]}. Therefore E is a splitting extension of the family \((f_{i})_{i \in I}\) .

PROPOSITION 5. --- Let K be a field, \((f_i)_{i \in I}\) a family of non-constant polynomials in K{[}X{]}, E an extension of K, and F, F' subextensions of E which are each a splitting extension of \((f_i)_{i \in I}\) . Then F = F'.

Let \(\mathbf{R}_i\) be the set of roots of \(f_i\) in \(\mathbf{E}\) and \(\mathbf{R} = \bigcup_{i\in I}\mathbf{R}_i\) . Since \(f_i\) is a product of

polynomials of the first degree lying in F{[}X{]}, we have \(R_{i} \subset F\) . By Def. 2, we have \(F = K(R)\) ; in the same way we find that \(F' = K(R)\) .

COROLLARY. --- Let K be a field, \((f_{i})_{i\in I}\) a family of non-constant polynomials in K{[}X{]} and F, F' splitting extensions of \((f_{i})_{i\in I}\) . Then there exists a K-isomorphism of F onto F'.

This follows from Prop. 5 and V, p.~13, Cor. of Prop. 4.

